% ============================================
% ===== CAPÍTULO 3: ITERACIONES =====
% ============================================
\newpage
\chapter{Iteraciones}\label{chap:iteraciones}
\section{Primera Iteración}
\subsection{Planificación y Análisis}
Siendo esta la primera iteración, es en la que más se desarrolla para que en las siguientes iteraciones todo resulte más sencillo, como ya se había planeado adjudicandole más horas que al resto de iteraciones.
Esta iteración se basa en 4 puntos: la creación de la función para detectar los \textit{items} continuables;
las configuraciones iniciales en CouchDB; la creación de las primeras alarmas --- después de crearlas, solo habrá que modificar su contenido
en el resto de iteraciones --- y la creación del contrato.


\vspace{0.6cm}

Uno de los objetivos fundamentales de esta iteración es sentar las bases del proyecto y definir las tecnologías que se emplearán. Esto abarca tanto el lenguaje de programación como el sistema de control de versiones. En este caso, ambas decisiones son heredadas: se usa \textit{JavaScript} --- cumpliendo así el requisto RNF05 y el RNF09 --- porque es el lenguaje en el que están desarrolladas las librerías internas del equipo de Billing de IONOS Cloud, y para el versionado se utiliza GitLab de manera habitual. Además, por seguridad, el código se respalda periódicamente tanto en GitHub como en un disco duro físico, como se explicó en la sección \hyperref[sec:plan-contingencias]{Plan de contingencias y riesgos}.


\vspace{0.6cm}

Referido a la detección de los \textit{items} continuables, se realiza en dos fases: la primera consiste en detectar solamente un SKU concreto y, cuando funcione, se extiende al resto de SKUs.

\vspace{0.6cm}

Tanto la configuración inicial de CouchDB, como la creación de las alarmas y la formalización del contrato se abordan con especial atención al detalle. En el caso del contrato, esto implica verificar que todos los SKUs disponibles estén contemplados, ya que una cobertura completa desde el inicio evita tener que revisarlo en iteraciones posteriores. Además, tanto el contrato como las alarmas se diseñan siguiendo el principio de abierto/cerrado de los principios SOLID: preparados para incorporar futuras actualizaciones mediante la adición de nuevos elementos, sin necesidad de modificar el código existente.

\vspace{1cm}

\subsubsection{Verificación de SKUs con PostMan}
Para la verificación de los SKUs disponibles en la API de IONOS Cloud, hemos optado por utilizar
\textbf{PostMan}\Cite{PostMan-Learning} como herramienta de prueba y automatización. Antes de comenzar a utilizarlo, es
importante comprender su funcionamiento y las capacidades que ofrece para este proyecto.

PostMan es una plataforma de colaboración para el desarrollo de APIs que permite diseñar, probar,
documentar y monitorizar APIs de manera eficiente. Para este trabajo, aprovecharemos especialmente
su capacidad para ejecutar scripts de prueba en JavaScript y su sistema de colecciones que permite
automatizar peticiones.

\paragraph{Conceptos básicos de PostMan para la verificación de SKUs}

Antes de comenzar a usarlo, es importante entender algunos conceptos
clave de PostMan:

\begin{itemize}
    \item \textbf{Collections:} Son grupos de peticiones relacionadas que pueden organizarse en carpetas.
          En nuestro caso, crearemos una colección específica para las pruebas de verificación de SKUs de
          IONOS Cloud.

    \item \textbf{Environments y variables:} PostMan permite definir entornos con variables reutilizables.
          Utilizaremos variables de entorno para almacenar:
          \begin{itemize}
              \item La URL base de la API de IONOS Cloud.
              \item Los tokens de autenticación de forma segura.
              \item Los resultados de las verificaciones para su uso en otras peticiones.
          \end{itemize}

    \item \textbf{Pre-request scripts:} Scripts que se ejecutan \textit{antes} de enviar una petición.
          Nos permitirán preparar datos, generar tokens temporales o configurar cabeceras de autenticación
          automáticamente.

    \item \textbf{Tests scripts:} Scripts que se ejecutan \textit{después} de recibir la respuesta de la API.
          Aquí implementaremos la lógica de verificación de SKUs, comparando los recursos obtenidos con
          la lista predefinida de recursos requeridos.

    \item \textbf{Console:} Herramienta de depuración que muestra los logs generados por nuestros scripts.
          Utilizaremos \texttt{console.log()} para generar reportes detallados del estado de cada verificación.

    \item \textbf{Runner:} Funcionalidad que permite ejecutar colecciones completas de manera automática,
          ideal para realizar verificaciones periódicas sin intervención manual.
\end{itemize}

\subsubsection{Base de datos con CouchDB}
Para las configuraciones, exceptuando los tokens y las claves, se ha optado por utilizar \textbf{CouchDB}\Cite{CouchDB-Info}, para cumplir el requisito RNF10, una base de datos NoSQL
orientada a documentos que destaca por su facilidad de integración con aplicaciones web. Realizando esta tarea el 20 de febrero de 2026.

\paragraph{Características principales de CouchDB para el proyecto}

\begin{itemize}
    \item \textbf{Almacenamiento basado en JSON:} Los datos se guardan como documentos JSON, lo que
          facilita su manipulación desde aplicaciones JavaScript como la que estamos desarrollando.

    \item \textbf{API RESTful nativa:} CouchDB expone una API HTTP que permite realizar operaciones
          CRUD mediante peticiones estándar (GET, POST, PUT, DELETE), lo que simplifica tanto la comunicación con
          el backend como la integración con nuestro proyecto.

    \item \textbf{Replicación maestro-maestro:} Soporta sincronización bidireccional entre múltiples
          instancias, lo que facilita la distribución de datos y la tolerancia a fallos, aspecto relevante para cumplir el requisito RNF01.

    \item \textbf{MVCC (Multiversion Concurrency Control):} Utiliza control de concurrencia mediante
          versiones, evitando bloqueos en lecturas y escrituras simultáneas, ayudando a cumplir el requisito RNF01.
\end{itemize}




\subsection{Desarrollo y despliegue}

\subsubsection{Creación del contrato}
El requisito RF01.01 exige un usuario exclusivo para las pruebas, por lo que es necesario crear un contrato dedicado. Tras una breve investigación se concluyó que no existía un método sencillo para crearlo de forma automática; se optó por hacerlo manualmente, ya que el coste de tiempo era menor y el proceso más directo. La creación del contrato se inició el 15 de febrero y finalizó el 19 de febrero de 2026.

El método elegido combina el DCD de IONOS con comprobaciones periódicas en PostMan. Como paso previo hay que disponer de una cuenta de IONOS Cloud y darse de alta en el DCD; al tratarse de un proceso sencillo, no se detalla aquí ---si se desea más información, se puede consultar el \hyperref[anexo:crear-cuenta-token]{Anexo 2}---.

Al añadir los \textit{items} de la lista al contrato se completa el requisito RF02.01 y sus dependencias. Una vez incorporados, estos recursos permanecen activos todos los meses, lo que satisface ---para los \textit{items} continuables--- el requisito RF02.03: sin modificaciones, el consumo mensual será predecible y estable.

Previamente a la creación del contrato, se elaboró la lista de SKUs a detectar a partir de la documentación de IONOS Cloud~\cite{billing-IONOSCloud}. Inicialmente constaba de 52 SKUs ---con posibles incorporaciones posteriores en las iteraciones 4 y 5, al añadir los \textit{items} no continuables---. El resultado se recoge en el \hyperref[anexo:Lista-de-SKUs]{Anexo 1}, con los SKUs seleccionados teniendo en cuenta los requisitos RF01.03 y RNF04, y dando cumplimiento al requisito RF01.02. Dicha lista se completó el 6 de febrero de 2026, respetando las fechas establecidas.

\vspace{0.5cm}

El primer problema encontrado fue vincular los recursos a sus SKUs, ya que la denominación de algunos no es intuitiva. Para resolverlo se recurrió a la API interna de IONOS Cloud~\cite{billing-IONOSCloud}.

Para utilizar dicha API se necesitaba, además de la documentación, un token generado por el contrato ---la forma de obtenerlo es sencilla y no es objeto del proyecto; para más información véase el \hyperref[anexo:crear-cuenta-token]{Anexo 2}---. Una vez obtenido, resulta imprescindible crear un plan de uso y control de estos tokens, dado que constituyen información delicada.

\subsubsection{Plan de gestión de tokens de autenticación}
\label{iteraciones:gestion-tokens}
Para garantizar la seguridad en el acceso a la API de IONOS Cloud, se establece un plan de gestión de tokens de autenticación utilizando \textbf{KeePassXC} como gestor de contraseñas especializado, cumpliendo así el requisito RNF03.

\paragraph{Justificación de la herramienta}
\vspace{0.25cm} 
KeePassXC es un gestor de contraseñas de código abierto, auditado de forma independiente, que proporciona cifrado robusto (AES-256, Twofish o ChaCha20) para la información almacenada \cite{keepassxc-audit}. Entre sus características más relevantes para este proyecto destacan:
\begin{itemize}
    \item Capacidad de generar contraseñas y tokens seguros.
    \item Almacenamiento cifrado de secretos con protección en memoria.
    \item Generación integrada de TOTP (Time-based One-Time Passwords) \cite{keepassxc-features}.
    \item Compatibilidad con archivos clave y YubiKey para autenticación de doble factor.
    \item Integración con navegadores mediante extensiones oficiales.
\end{itemize}

\paragraph{Ciclo de vida de los tokens en IONOS Cloud}
\vspace{0.25cm}
Los tokens de autenticación en IONOS Cloud presentan las siguientes características \cite{ionos-token-management}:
\begin{itemize}
    \item Se generan a través del \textit{Authentication Token Manager} en el DCD.
    \item Cada token tiene un TTL (\textit{Time To Live}) configurable entre 1 hora y 365 días.
    \item El valor del token solo se muestra una única vez durante su generación.
    \item Es posible descargar el token como archivo de texto en el momento de creación.
    \item Se pueden generar hasta 100 tokens por cuenta.
    \item La eliminación de un token lo invalida inmediatamente, incluso antes de su fecha de expiración.
\end{itemize}

\paragraph{Procedimiento de gestión segura}
\vspace{0.25cm}
\begin{enumerate}
    \item \textbf{Generación del token:}
          \begin{itemize}
              \item Acceder al DCD: Menú → Management → Token Manager.
              \item Seleccionar TTL adecuado al caso de uso (para desarrollo: 30-90 días; para producción: evaluar periodicidad de rotación).
              \item Generar el token y, en la pantalla de visualización, seleccionar \textit{Download} para obtener el archivo de texto.
          \end{itemize}

    \item \textbf{Almacenamiento en KeePassXC:}
          \begin{itemize}
              \item Crear una entrada específica para cada token en la base de datos de KeePassXC.
              \item Campos a rellenar:
                    \begin{itemize}
                        \item \textbf{Título:} ``Token IONOS API - [Propósito]'' (ej: ``Token IONOS API - Desarrollo Iteración 1'').
                        \item \textbf{Nombre de usuario:} Identificador del token (Token ID).
                        \item \textbf{Contraseña:} Pegar el valor del token.
                        \item \textbf{URL:} https://api.ionos.com (para integración con navegador).
                        \item \textbf{Notas:} Fecha de generación, TTL seleccionado, fecha de expiración calculada, propósito específico.
                    \end{itemize}
              \item Adjuntar el archivo de texto descargado como archivo adjunto en la entrada.
              \item Configurar la entrada para generar TOTP si se requiere autenticación adicional.
          \end{itemize}

    \item \textbf{Rotación y expiración:}
          \begin{itemize}
              \item Configurar un recordatorio en el calendario para 7 días antes de la fecha de expiración.
              \item En la fecha de rotación:
                    \begin{enumerate}
                        \item Generar nuevo token en el DCD.
                        \item Almacenarlo en KeePassXC siguiendo el procedimiento anterior.
                        \item Eliminar el token antiguo desde el DCD.
                        \item Marcar la entrada antigua como expirada o eliminarla de la base de datos.
                    \end{enumerate}
          \end{itemize}

    \item \textbf{Seguridad adicional:}
          \begin{itemize}
              \item La base de datos de KeePassXC debe protegerse con una contraseña maestra robusta (mínimo 15 caracteres con alta entropía).
              \item Habilitar el bloqueo automático de la base de datos tras periodo de inactividad.
          \end{itemize}
\end{enumerate}

\subsubsection{Uso del Token}

Una vez definida la gestión de los tokens, estos se pueden utilizar para realizar llamadas a la API de facturación. En concreto, se emplea el endpoint \textit{"https://api.ionos.com/billing/\{contract\}/products"}~\cite{billing-IONOSCloud}, que devuelve la información de los productos asociados al contrato. De este modo es posible verificar la presencia de cada recurso deseado.

Para ejecutar estas llamadas se utiliza PostMan.
Como primer SKU a detectar se opta por A01000 (IP estática). Para ello se añade una IP estática de forma manual en el DCD y, tras pasar el periodo que tarda en actualizarse
la base de datos de la API, se hace una llamada y debe aparecer. En este caso debe salir algo similar a:

\begin{lstlisting}[caption=Ejemplo de producto en API de IONOS, label=lst:producto-json]
{
    "products": [
        {
            "meterId": "A01000",
            "meterDesc": "30d per static IP address",
            "deprecated": false,
            "unit": "30Days",
            "type": "consumption",
            "unitCost": {
                "quantity": "5",
                "unit": "EUR"
            }
        }
    ]
}
\end{lstlisting}

Realizar este proceso de forma manual ---tanto la verificación como la incorporación de productos--- resulta propenso a errores por su naturaleza repetitiva. Por ello, se opta por una solución más fiable y reutilizable: aprovechar la capacidad de PostMan para ejecutar scripts y crear uno que compruebe automáticamente si están presentes los SKUs deseados.

\subsubsection{Creación del script de verificación}
\label{iteraciones:script-verificacion-postman}

Se creó un script en JavaScript para PostMan que verifica los SKUs disponibles.

\paragraph{Funcionamiento}
\begin{enumerate}
    \item Define un objeto \texttt{requiredResources} con 14 categorías de SKUs (IPs estáticas, VMs, almacenamiento, bases de datos, etc.).
    \item La función \texttt{checkResources()} compara cada SKU requerido con los devueltos por la API.
          \begin{lstlisting}[label=lst:verificacion-categorias]
Object.keys(requiredResources).forEach(category => {
    const resources = requiredResources[category];
    results.available[category] = [];
    results.missing[category] = [];
    
    resources.forEach(resource => {
        const found = products.find(p => p.meterId === resource);
        if (found) { results.available[category].push(resource);
        } else { results.missing[category].push(resource); } }); });

\end{lstlisting}
    \item Clasifica los resultados en \texttt{available} (encontrados) y \texttt{missing} (no encontrados).
          \begin{figure}[h!]
              \centering
              \includegraphics[width=1\textwidth]{PostMan-Test.png}
              \caption{Prueba Test}
              \label{fig:edt}
          \end{figure}
    \item Detecta patrones especiales: VMs CUBE (contienen ``CUBES'') y licencias Windows (empiezan por ``WL'').
    \item Almacena los resultados en variables de entorno de PostMan.
    \item Genera un reporte en consola con el estado de cada categoría.
    \item Calcula estadísticas globales: total disponibles vs total requeridos.
          \begin{figure}[h!]
              \centering
              \includegraphics[width=0.25\textwidth]{PostMan-Resumen.png}
              \caption{Prueba Resumen Test}
              \label{fig:edt}
          \end{figure}
    \item Ejecuta tests automáticos (código 200, productos no vacíos).
          \begin{figure}[h!]
              \centering
              \includegraphics[width=0.7\textwidth]{PostMan-CodsOK.png}
              \caption{Prueba Cods. Test}
              \label{fig:edt}
          \end{figure}
\end{enumerate}

\paragraph{Ventajas}

Automatiza la verificación, proporciona trazabilidad, se integra en PostMan y es fácilmente mantenible. Además, esta estructura es migrable al proyecto, permitiendo tener una parte ya hecha.

\subsection{Integración del código}
\subsubsection{Configuración de Entornos}
Al ser la primera iteración se deben configurar los entornos que se van a usar.
Esto consiste en la creación del proyecto, su vinculación a GitLab y GitHub,
el plan de copias de seguridad en físico y la configuración inicial en CouchDB. Estas tareas son indispensables para cumplir los requisitos RNF10, RF02.02, RF04.03 y RF05.01.  
\subsubsection{Despliegue del proyecto y vinculación a GitLab}
El proyecto se crea directamente desde GitLab ---cumpliendo así el requisito RNF08--- porque se desea integrarlo en el grupo de trabajo del equipo de \textit{Cloud Billing} y \textit{ERP} de IONOS Cloud, de modo que, tras la finalización del proyecto, todo el equipo pueda continuar con su mantenimiento. Por este motivo se sigue, además, la estructura de proyectos propia del equipo. La única peculiaridad fue la necesidad de usar la VPN, ya que el GitLab empleado es el corporativo del equipo, alojado en Alemania.

\subsubsection{Estrategia de copias de seguridad del código fuente usando GitHub y disco físico}
Para garantizar la integridad y disponibilidad del código fuente desarrollado en este proyecto,
se ha establecido una estrategia de copias de seguridad multicapa que complementa el sistema de
control de versiones principal.

\paragraph{Uso de GitLab como sistema de control de versiones principal}
GitLab actúa como el repositorio principal donde se gestiona el desarrollo activo del proyecto.
Sus características incluyen:
\begin{itemize}
    \item Control de versiones con historial completo de cambios.
    \item Gestión de ramas para el desarrollo de funcionalidades.
    \item Integración continua para pruebas automáticas.
    \item Documentación integrada mediante wikis.
\end{itemize}

\paragraph{GitHub como repositorio espejo de respaldo}
A pesar de utilizar GitLab como plataforma principal, se ha configurado GitHub como repositorio
secundario para proporcionar una capa adicional de redundancia. Esta decisión se fundamenta en:

\begin{itemize}
    \item \textbf{Independencia entre plataformas:} Al estar alojados en servidores de diferentes
          proveedores, una caída o problema en GitLab no compromete el acceso al código.

    \item \textbf{Redundancia geográfica:} Los datos se replican en centros de datos de diferentes
          ubicaciones, protegiendo frente a desastres regionales. Esto previene uno de los posibles riesgos estudiados en el plan de riesgos.
\end{itemize}

\paragraph{Copia de seguridad en disco físico local}
Complementando las copias en la nube, se realiza una copia de seguridad periódica en un disco físico
externo siguiendo estos principios: 

\begin{itemize}
    \item \textbf{Frecuencia de respaldo:} Copia completa del repositorio cada 15 días y copia
          incremental de los cambios diarios.

    \item \textbf{Almacenamiento seguro:} El disco se mantiene en una ubicación física diferente a la
          del equipo de desarrollo, protegiendo frente a robo, incendio o fallo del equipo principal.

    \item \textbf{Verificación de integridad:} Tras cada backup, se ejecuta una comprobación automática
          mediante \texttt{git fsck} para garantizar que el repositorio no esté corrupto.
\end{itemize}

\subsubsection{Configuración en CouchDB}
Debido a que el proyecto está integrado en el grupo de trabajo de \textit{Cloud Billing} y \textit{ERP}, se utiliza su grupo de bases de datos. --- Cumpliendo así el requisito RNF10. ---
Lo único que hay que tener en cuenta es el uso de las credenciales; dado que las credenciales para acceder a CouchDB
son delicadas ---ya que se accede a todas las configuraciones del equipo---, se emplea un protocolo similar al descrito en la
sección \hyperref[iteraciones:gestion-tokens]{Plan de gestión de tokens de autenticación}.

Una vez creada la base de datos propia, se crea un documento para vincularlo con el proyecto.


\paragraph{Estructura del documento de configuración}
El documento almacenado en CouchDB contiene la lógica y parámetros necesarios para las verificaciones:

\begin{itemize}
    \item \textbf{Identificador:} \texttt{billing-reference-verifier} - Documento raíz que centraliza
          toda la configuración del verificador. Este nombre coincide con el del proyecto, estandarizado por el grupo de trabajo, lo que cumple parcialmente el requisito RNF09.

    \item \textbf{Programación de tareas (scheduler):} Define ejecuciones automáticas periódicas; sobre esto se desarrollará en mayor profundidad cuando se hable de ello en concreto.

\end{itemize}

\paragraph{Configuración específica para verificación de SKUs}
El módulo \texttt{utilization} contiene la lista de SKUs a monitorizar, siguiendo la siguiente estructura:

\begin{lstlisting}[caption=Almacenamiento SKUs en CouchDB, label=lst:producto-json]
{
    "skus_monitorizados": [
        "A01000", "C02000", "C03000", "C04000", "C05000", "C06000",
        [...],
        "CWSQL1001", "RHEL2100"
    ]
}
\end{lstlisting}
\subsubsection{Estructura general del proyecto billing-reference-verifier}

El proyecto \texttt{billing-reference-verifier} es una API REST interna desarrollada para IONOS
cuyo propósito principal es verificar la utilización de recursos y validar facturas asociadas a
contratos de facturación. Antes de abordar los detalles técnicos, se describe a continuación la estructura general del proyecto y las tecnologías que lo componen. Cumpliendo asi los requisitos RNF05, RNF06, RNF07 y  RNF09.

\paragraph{Información general}
\begin{itemize}
    \item \textbf{Nombre:} billing-reference-verifier
    \item \textbf{Tecnología:} Node.js (versión $\geq$ 18.0.0) con framework interno \texttt{@internal/core}
    \item \textbf{Tipo:} API REST interna
\end{itemize}

\paragraph{Arquitectura de directorios}
La estructura del proyecto se organiza de la siguiente manera:

\begin{verbatim}
billing-reference-verifier/
|-- package.json                  # Dependencias y scripts de Node.js
|-- README.md                     # Documentacion principal
|-- run.sh / start.sh             # Scripts de inicio del servicio
|-- config/
|   +-- nginx.conf                # Configuracion de proxy reverso
|-- docs/
|   |-- README.md                  # Documentacion adicional
|   +-- SUMMARY.md                 # Resumen de documentacion
|-- src/
|   |-- run.js                     # Punto de entrada (bootstrap)
|   |-- helpers/                    # Clases auxiliares y utilidades
|   |   |-- actions.js              # Constantes de eventos/acciones
|   |   |-- invoice.js              # Helper para operaciones de facturas
|   |   |-- schema.js               # Esquemas de validacion Joi
|   |   +-- utilization.js          # Helper para datos de utilizacion
|   +-- modules/                    # Modulos funcionales de la aplicacion
|       |-- utilization/
|           +-- index.js            # Logica de verificacion de utilizacion
+-- test/
    +-- test-usage-validate.js      # Pruebas automatizadas
\end{verbatim}

\paragraph{Componentes principales}

\paragraph{Punto de entrada (run.js)}
Realiza el bootstrap de la aplicación utilizando el framework \texttt{@internal/core}, carga la
configuración definida en \texttt{package.json} e inicializa el servidor y los módulos.

\paragraph{Módulos funcionales}
\begin{itemize}
    \item \textbf{utilization:} Verifica las cantidades diarias de uso frente a los valores
          esperados para los SKUs configurados.
\end{itemize}

\paragraph{Helpers}
\begin{itemize}
    \item \textbf{actions.js:} Define constantes para los eventos del mediador (como
          \texttt{APPLICATION\_RUN}, \texttt{UTILIZATION\_VERIFIER\_COMPLETE}, etc.).

    \item \textbf{schema.js:} Contiene esquemas Joi para validar las respuestas de la API
          (\texttt{invoiceSchema}, \texttt{utilizationSchema}, \texttt{usageSchema}).

    \item \textbf{utilization.js:} Implementa la clase \texttt{UtilizationHelper} para obtener
          datos de los endpoints \texttt{/utilization} y \texttt{/usage} de la API de facturación.
\end{itemize}

\paragraph{Endpoints REST disponibles}

\begin{itemize}
    \item \texttt{GET /data/v1/utilization/check/:contractId} --- Verifica la presencia de ítems configurados.
    \item \texttt{GET /data/v1/utilization/debug} --- Información de diagnóstico.
    \item \texttt{POST /data/v1/invoice/manual-test/:contractId} --- Test manual de validación.
\end{itemize}

\paragraph{Flujo de datos}

\begin{enumerate}
    \item Una petición HTTP llega al módulo \texttt{data-api}, que valida los headers (incluyendo
          el token Bearer).
    \item El mediador emite un evento al módulo correspondiente según el endpoint solicitado.
    \item El módulo utiliza los helpers para realizar llamadas a la \texttt{billing-api} de IONOS.
    \item Las respuestas se validan mediante los esquemas Joi definidos en \texttt{schema.js}.
    \item Se devuelve una respuesta JSON estructurada con el resultado, resumen y timestamp.
\end{enumerate}

\paragraph{Dependencias externas}
\begin{itemize}
    \item \textbf{@internal/core:} Framework base que proporciona servidor Koa, cliente REST y utilidades.
    \item \textbf{helmet:} Middleware de seguridad para cabeceras HTTP.
    \item \textbf{@internal/dev:} Herramientas de desarrollo (eslint, nodemon, etc.).
    \item \textbf{joi:} Validación de esquemas (integrado a través de @internal/core).
\end{itemize}

\paragraph{Patrones de diseño utilizados}
\begin{itemize}
    \item \textbf{Módulos desacoplados:} Cada módulo hereda de una clase \texttt{Module} base.
    \item \textbf{Mediator Pattern:} Comunicación entre módulos mediante eventos.
    \item \textbf{Factory Pattern:} Uso de \texttt{createHelper()} para instanciar helpers.
    \item \textbf{Dependency Injection:} Container para servicios.
    \item \textbf{Schema Validation:} Joi para validar estructuras de datos.
\end{itemize}

\paragraph{Autenticación }
Todas las peticiones requieren el header \texttt{Authorization: Bearer <token>}. Este token se
transmite a la \texttt{billing-api} de IONOS para autenticar las operaciones y se gestiona como ya se explicó en
\hyperref[iteraciones:gestion-tokens]{el plan de gestión de tokens de autenticación} y garantizar el cumplimiento del requisito RNF03.

\subsection{Implementación de función de Detección de items continuables}
Una vez definida la estructura del proyecto, el 21 de febrero se inicia la implementación de la función de detección de ítems continuables en el módulo \texttt{utilization}. Esta función realiza llamadas a la API de facturación para obtener los datos de utilización y compararlos con la lista de SKUs monitorizados en CouchDB, logrando así la primera parte del requisito RNF04.02.

Para ello se implementa la función \texttt{checkUtilization()}, basada en el script que se creó para PostMan ---explicado en la sección \hyperref[iteraciones:script-verificacion-postman]{Creación del script de verificación}---. Al igual que en PostMan, el desarrollo se realiza en dos fases: primero se detecta un SKU concreto y, una vez verificado su correcto funcionamiento, se extiende al resto de SKUs.

\paragraph{Validación de respuestas con Joi}

Para garantizar la integridad de los datos recibidos de la API de facturación, se ha implementado la validación de las respuestas mediante esquemas \texttt{Joi}, definidos en el helper \texttt{schema.js}. Esto permite detectar de forma temprana respuestas inesperadas o malformadas de la API, devolviendo un error controlado en lugar de propagar datos incorrectos al resultado de la verificación. Además se ha creado la verificación de todos los futuros datos necesarios, es decir, ahora solo se necesitan los datos de utilización, pero se ha creado la validación de las facturas y del uso para futuras iteraciones, lo que cumple el requisito RNF06.

El archivo \texttt{schema.js} define tres esquemas de validación, cada uno adaptado a un endpoint distinto de la API de facturación:

\begin{itemize}
    \item \textbf{\texttt{invoiceSchema}:} Valida las respuestas de facturas procedentes de NetSuite. Incluye la verificación del total facturado (unidad y cantidad), metadatos de la factura (fechas, identificadores de factura, contrato y cliente) y un array de \textit{datacenters}, cada uno con su ubicación, medidores y descuentos (\textit{rebates}). Un aspecto destacable es la validación del formato del identificador de factura mediante una expresión regular:
\end{itemize}

\begin{lstlisting}[caption=Validación del formato de invoiceId, label=lst:invoice-pattern]
metadata: joi.object({
    invoiceId: joi.string().pattern(/^GY\d+$/).required(),
    contractId: joi.string().required(),
    postingPeriod: joi.string().pattern(/^\d{4}-\d{2}$/).required()
    // ...
}).required()
\end{lstlisting}

Esto garantiza que el identificador de factura siga siempre el patrón esperado por el sistema, rechazando respuestas con formatos incorrectos. De forma similar, el \texttt{postingPeriod} se valida contra el formato \texttt{YYYY-MM}.

\begin{itemize}
    \item \textbf{\texttt{utilizationSchema}:} Valida las respuestas del endpoint \texttt{/utilization}, que es el utilizado en esta iteración para detectar los ítems continuables. La estructura es más compleja, ya que cada medidor incluye referencias UUID a recursos y servidores, tipo de recurso, rango de fechas y región geográfica. Un ejemplo representativo de la validación de un medidor:
\end{itemize}

\begin{lstlisting}[caption=Validación de un medidor en utilizationSchema, label=lst:utilization-meter]
meters: joi.array().items(
    joi.object({
        id: joi.string().uuid().required(),
        resourceId: joi.string().uuid().required().allow(null),
        serverId: joi.string().uuid().required().allow(null).allow(''),
        type: joi.string().required(),
        meterId: joi.string().required(),
        region: joi.string().required(),
        quantity: joi.object({
            quantity: joi.number().required(),
            unit: joi.string().required()
        }).required()
    }).required()
).required()
\end{lstlisting}

Es relevante el uso de \texttt{.allow(null)} y \texttt{.allow('')} en campos como \texttt{resourceId} y \texttt{serverId}, ya que no todos los recursos del DCD tienen un servidor o recurso físico asociado de forma directa, pero el campo debe estar presente en la respuesta. Sin esta flexibilidad, la validación rechazaría respuestas legítimas de la API.

\begin{itemize}
    \item \textbf{\texttt{usageSchema}:} Valida las respuestas del endpoint \texttt{/usage}, que se utilizará en iteraciones posteriores. Su estructura es más sencilla que la del esquema de utilización, pero incorpora un aspecto interesante: la cantidad puede venir como número o como cadena numérica, lo que se resuelve con \texttt{alternatives()}:
\end{itemize}

\begin{lstlisting}[caption=Validación flexible de cantidad en usageSchema, label=lst:usage-alternatives]
quantity: joi.object({
    quantity: joi.alternatives().try(
        joi.number(),
        joi.string().regex(/^\d+(?:\.\d+)?$/)
    ).required(),
    unit: joi.string().required()
}).required()
\end{lstlisting}

Esta flexibilidad es necesaria porque la API de facturación devuelve en ocasiones valores numéricos como cadenas de texto (por ejemplo, \texttt{"4.5"} en lugar de \texttt{4.5}), comportamiento que puede variar entre versiones o entre distintos tipos de recursos.

Los tres esquemas comparten un conjunto de validadores comunes de Joi: \texttt{required()} para campos obligatorios, \texttt{optional()} para opcionales, \texttt{pattern()} para expresiones regulares, \texttt{valid()} para valores permitidos (como \texttt{'EU'} y \texttt{'US'} en la ubicación de los datacenters), \texttt{uuid()} para identificadores universales y \texttt{array().items()} para validar arrays de objetos con estructura definida.

\subsubsection{Creación de endpoint y decisión de formato de URL}
Debido a que no hay ningún requisito que especifique el formato de la URL, se opta por un formato RESTful que refleje claramente la acción realizada. Se decide utilizar el endpoint \texttt{GET /data/v1/utilization/check/:contractId}, donde \texttt{contractId} es un parámetro dinámico que permite especificar el contrato para el cual se desea realizar la verificación de utilización. Este formato es intuitivo y sigue las convenciones REST, facilitando su comprensión y uso por parte de otros desarrolladores o sistemas que puedan integrarse con esta API en el futuro.
Sin embargo, no se puede garantizar que este formato se mantenga en el futuro, ya que no hay un requisito que lo especifique, por lo que podría cambiarse en futuras iteraciones si el cliente lo considera necesario. Por ello, se documenta claramente en el README del proyecto para asegurar que cualquier cambio futuro se realice de manera controlada y con conocimiento de las partes involucradas.


\subsection{Desarrollo e implementación}

La primera tarea llevada a cabo en esta primera iteración ha sido la configuración de los entornos necesarios para el desarrollo del sistema. Esta fase incluye principalmente el despliegue del proyecto en GitLab, la configuración de la base de datos CouchDB y su vinculación con el proyecto, y la preparación del entorno de pruebas mediante PostMan. Aunque también se han creado los repositorios de respaldo en GitHub y se ha configurado la VPN para acceder a la infraestructura corporativa, estas acciones se consideran triviales dentro del alcance técnico del proyecto.

\subsubsection{Configuración de PostMan}
Para la realización de pruebas sobre las peticiones HTTP a la API de facturación de IONOS Cloud, se ha configurado PostMan. La herramienta permite simular las llamadas que realizará el sistema de verificación, facilitando la creación y prueba de peticiones REST antes de integrarlas en el código.

Se han configurado los siguientes elementos:
\begin{itemize}
    \item \textbf{Collection:} Se ha creado una colección específica denominada \textit{billing-reference-verifier} para agrupar todas las peticiones de verificación.
    \item \textbf{Variables de entorno:} Se han definido variables para la URL base de la API de IONOS Cloud, el \texttt{contractId} de pruebas y el token de autenticación Bearer, evitando así tener que introducirlos manualmente en cada petición.
\end{itemize}

Es importante configurar correctamente los encabezados, el cuerpo de las peticiones según lo requiera cada endpoint de la API de facturación, ya que una configuración incorrecta puede provocar errores difíciles de diagnosticar.

Tras reunirse con el cliente, se sugerieron por su parte dos cambios, la primera el estandarizar el contrato y eliminarlo de la URL de la petición, denido a que siempre se usará el mismo contrato y en el caso de querer usar otro en el futuro se puede modificar en la configuración. La segunda fue usar desde PostMan \texttt{Authorization: Bearer <token>} en lugar de una variable personalizada, para que sea más fácil de usar y entender. Ambos cambios se implementarán en la proxima iteración.

%% TODO: Si quieres incluir el esquema Joi utilizado (utilizationSchema, usageSchema, etc.), añádelo aquí.


\subsubsection{Configuración de ejecución programada}

Para automatizar las verificaciones y cumplir el requisito RF04.01, se ha configurado un \textit{scheduler} en el documento de CouchDB que define la periodicidad de ejecución de las comprobaciones.

%% TODO: Añadir aquí los detalles concretos del scheduler: frecuencia (diaria, semanal, etc.),
%% formato de la configuración en CouchDB y cómo se dispara la ejecución (cron job, servicio interno, etc.).

De esta forma, el sistema realiza verificaciones de manera autónoma sin necesidad de intervención manual, generando alertas únicamente cuando se detectan discrepancias.

\subsection{Pruebas y validación}

Para verificar el correcto funcionamiento de la iteración se han diseñado pruebas a dos niveles: pruebas manuales con PostMan y pruebas automatizadas integradas en el proyecto.

\subsubsection{Pruebas manuales con PostMan}

Se han ejecutado peticiones al endpoint \texttt{GET /data/v1/utilization/check/:contractId} desde PostMan, verificando que:
\begin{itemize}
    \item La respuesta devuelve un código 200 cuando la petición es correcta.
    \item Todos los SKUs de la lista se encuentran en la categoría \texttt{available}.
    \item La respuesta incluye la estructura esperada: \texttt{result}, \texttt{summary} y \texttt{timestamp}.
    \item Al eliminar un recurso del contrato de pruebas de forma temporal, el SKU correspondiente aparece en la categoría \texttt{missing} y se dispara la alerta configurada.
\end{itemize}


\subsubsection{Pruebas automatizadas}

Se han implementado pruebas automatizadas en el archivo \texttt{test/test-usage-validate.js} que verifican la correcta validación de los esquemas Joi frente a las respuestas de la API. El enfoque elegido consiste en definir dos muestras de datos ---una válida y otra inválida--- y comprobar que el esquema acepta la primera y rechaza la segunda.

La muestra válida reproduce una respuesta real del endpoint \texttt{/usage}, con un datacenter de ubicación \texttt{EU}, un medidor con SKU \texttt{A01000} (IP estática) y una cantidad expresada como cadena numérica (\texttt{"0.87"}), lo que permite verificar simultáneamente la estructura general y el comportamiento del validador \texttt{alternatives()} descrito anteriormente:

\begin{lstlisting}[caption=Muestra válida para el test de validación, label=lst:valid-sample]
const validSample = {
    startDate: '2025-11-01',
    endDate: '2025-11-26',
    datacenters: [{
        id: 'a786c097-d6a6-531c-a55d-6e72ce4c35c5',
        name: '',
        location: 'EU',
        meters: [{
            meterId: 'A01000',
            meterDesc: '30d per static IP address',
            quantity: { quantity: '0.87', unit: '30Days' }
        }]
    }],
    metadata: { customerId: '117048209', contractId: '036348704' }
};
\end{lstlisting}

La muestra inválida introduce intencionadamente dos errores: una cantidad con un valor no numérico (\texttt{"not-a-number"}) y la ausencia del campo obligatorio \texttt{contractId} en los metadatos. De este modo se comprueba que el esquema detecta tanto errores de formato como campos ausentes:

\begin{lstlisting}[caption=Muestra inválida para el test de validación, label=lst:invalid-sample]
const invalidSample = {
    // ...misma estructura...
    meters: [{
        meterId: 'A01000',
        meterDesc: '30d per static IP address',
        quantity: { quantity: 'not-a-number', unit: '30Days' }
    }],
    metadata: { customerId: '117048209' } // contractId ausente
};
\end{lstlisting}

Las aserciones se realizan con el módulo estándar \texttt{assert} de Node.js, verificando que la validación de la muestra válida no produzca errores y que la muestra inválida sí los genere:

\begin{lstlisting}[caption=Aserciones del test, label=lst:test-assertions]
const { error: vError } = usageSchema.validate(validSample, { convert: true });
assert.strictEqual(vError, undefined, 'validSample should pass');

const { error: iError } = usageSchema.validate(invalidSample, { convert: true });
assert(iError, 'invalidSample should fail');
\end{lstlisting}

La opción \texttt{convert: true} permite que Joi transforme automáticamente los tipos cuando sea posible (por ejemplo, cadenas de fecha a objetos \texttt{Date}), simulando así el comportamiento real del sistema en producción. El test se ejecuta mediante \texttt{node test/test-usage-validate.js} y devuelve un código de salida 0 en caso de éxito o 1 en caso de fallo, lo que facilita su integración en pipelines CI/CD.


\subsection{Configuración de alertas}

Para cumplir los requisitos RF05.01 y RF05.02, se ha implementado un sistema de notificaciones que se dispara cuando la verificación detecta SKUs ausentes. En esta primera iteración, la alerta informa de los SKUs que no han sido encontrados en el contrato, permitiendo al equipo actuar de forma inmediata.

La configuración de alertas definitiva se realizará utilizando la infraestructura propia del equipo de Billing y ERP. Sin embargo, al tratarse de una primera iteración, se ha optado por una solución provisional más sencilla: el envío de notificaciones mediante \textit{webhooks} de Google Chat. Este mecanismo permite enviar mensajes a un espacio de Google Chat mediante una simple petición HTTP POST a una URL generada por la propia plataforma, sin necesidad de configurar servicios intermedios ni gestionar autenticación adicional.

El contenido de la alerta incluye:


Esta solución es suficiente para las primeras fases del proyecto, ya que permite al equipo recibir notificaciones en tiempo real sin un coste de implementación elevado. En la Figura~\ref{fig:alerta-iteracion1} se muestra un ejemplo de alerta recibida en Google Chat durante esta iteración.

\begin{figure}[H]
    \centering
    \includegraphics[width=0.8\textwidth]{Alerta1Iteracion.png}
    \caption{Ejemplo de alerta de la primera iteración recibida en Google Chat}
    \label{fig:alerta-iteracion1}
\end{figure}

En iteraciones posteriores, cuando se integre el sistema de alertas del equipo, bastará con sustituir la llamada al webhook por el mecanismo correspondiente, sin modificar la lógica de detección que genera el contenido de la alerta.

\subsection{Conclusiones de la iteración}

Tras finalizar esta primera iteración el 26 de febrero de 2026, me reuní con el cliente para revisar los resultados obtenidos y planificar los siguientes pasos. Se verificó que la función de detección de ítems continuables funciona correctamente, identificando tanto los SKUs presentes como los ausentes en el contrato de pruebas. Además, se confirmó que las alertas se generan y envían a Google Chat según lo esperado. Quedan pendientes para la próxima iteración los cambios sugeridos por el cliente en relación con las peticiones.

En cuanto a la planificación temporal, aunque la iteración se ha finalizado antes de lo previsto ---debido a los márgenes del plan de contingencias y a que se han invertido más horas semanales de las planificadas---, las horas totales dedicadas se han aproximado bastante a la estimación inicial. 

\begin{table}[H]
    \centering
    \begin{tabular}{|l|c|c|c|}
        \hline
        \textbf{Tarea} & \textbf{Estimadas (h)} & \textbf{Reales (h)} & \textbf{Desviación}\\ 
        \hline
        Detección de items continuables & 16h & 16h & 0\% \\
        Creación de contrato completo & 12h & 16h & +25\% \\
        Configuración de credenciales & 6h & 4h & -33.3\% \\
        Testeo de funcionamiento & 4h & 4h & 0\% \\
        Configuración de alertas y rutinas & 4h & 2h & -50\% \\
        \hline
        \textbf{Total} & \textbf{42h} & \textbf{42h} & \textbf{0\%} \\
        \hline
    \end{tabular}
    \caption{Comparativa de horas estimadas vs reales en la Iteración 1}
    \label{tab:horas-iteracion1}
\end{table}


\section{Iteración 2}
\subsection{Planificación}
En esta iteración se trabajará principalmente en ampliar las funcionalidades del sistema de verificación, evolucionando desde la simple detección de existencia de SKUs hacia la verificación de cantidades de uso. Además, se implementarán las mejoras solicitadas por el cliente durante la revisión de la primera iteración.

\textbf{Verificación de cantidades de uso:} En la iteración anterior se implementó la detección de existencia de SKUs (\texttt{existence-only}). Ahora se ampliará esta funcionalidad para verificar que las cantidades de uso reportadas por la API coincidan con los valores esperados. Esto permitirá detectar anomalías como recursos que reportan menos cores, RAM o almacenamiento del esperado. Además de ser el preambulo final para lograr el requisito RF04.02. 

\textbf{Sistema de thresholds:} Se implementará un sistema de umbrales configurable que permita distinguir entre situaciones de \textit{warning} (desviación menor) y \textit{error} (desviación crítica), adaptándose a las diferentes tolerancias según el tipo de recurso.

\textbf{Cambios solicitados por el cliente:} Se implementarán los dos cambios sugeridos en la reunión de cierre de la iteración anterior. Por un lado, estandarizar el contrato eliminándolo de la URL del endpoint ---ya que siempre se utilizará el mismo contrato de pruebas y, si se desea cambiar en el futuro, se modificará en la configuración de CouchDB---. Por otro lado, modificar PostMan para utilizar \texttt{Authorization: Bearer <token>} de forma estándar, mejorando la claridad y facilitando el uso por parte de otros miembros del equipo.


\subsection{Análisis} 

\subsubsection{Tipos de verificación: existence-only vs precise}
Tras analizar los diferentes SKUs monitorizados, se identifican dos categorías de verificación necesarias:

\begin{itemize}
    \item \textbf{existence-only:} Para recursos cuyo consumo puede variar legítimamente (IPs estaticas, Bases de datos). En estos casos, basta con verificar que existe consumo mayor que cero, sin comparar contra un valor esperado concreto.
    
    \item \textbf{precise:} Para recursos con cantidades fijas y predecibles (cores de CPU, GB de RAM). Estos requieren calcular el valor esperado y compararlo con el reportado por la API, aplicando umbrales de tolerancia para detectar discrepancias.
\end{itemize}

Por ejemplo, el SKU \texttt{C05000} (cores de CPU) tiene un valor base de 8 cores que debe mantenerse constante, mientras que \texttt{MKS1000} (almacenamiento en Kubernetes) solo requiere verificar que existe consumo.

\subsubsection{Sistema de thresholds}
Para los SKUs con verificación \texttt{precise}, se definen dos niveles de alerta basados en el ratio entre el valor real y el esperado:

\begin{itemize}
    \item \textbf{Warning:} Se activa cuando el ratio cae por debajo del umbral de warning (típicamente 0.99). Indica una desviación menor que requiere atención pero no es crítica.
    
    \item \textbf{Error:} Se activa cuando el ratio cae por debajo del umbral de error (típicamente 0.975). Indica una discrepancia significativa que requiere acción inmediata.
\end{itemize}

Estos umbrales son configurables por SKU en CouchDB, permitiendo ajustar la sensibilidad según las características de cada recurso. Ya que hay recursos como los cores de CPU donde una desviación del 1\% puede ser relevante, mientras que en otros como el almacenamiento en Kubernetes, la mera ausencia de consumo es suficiente para generar una alerta.

\subsubsection{Estandarización del contrato en la configuración}
Actualmente, el \texttt{contractId} se pasa como parámetro en la URL del endpoint (\texttt{/check/:contractId}). Tras analizar los casos de uso previstos, se concluye que el sistema operará siempre sobre un único contrato de referencia, por lo que tiene sentido centralizar este valor en la configuración del scheduler de CouchDB.

Es importante aclarar que el contrato no se elimina del endpoint en ningún momento. Sin embargo, en los endpoints de chequeo, test y debug (\texttt{/check}, \texttt{/test}, \texttt{/debug}) se omitirá el parámetro \texttt{contractId} de la URL, ya que estos endpoints utilizarán siempre el contrato de pruebas configurado en CouchDB. En cambio, para el resto de endpoints se mantiene el \texttt{contractId} como parámetro en la URL, permitiendo así modificar fácilmente el contrato objetivo sin necesidad de alterar la configuración almacenada en la base de datos.

\subsection{Desarrollo y pruebas}

\subsubsection{Configuración de itemsToCheck en CouchDB}
Se ha ampliado la estructura de configuración del módulo \texttt{utilization} en CouchDB para soportar ambos tipos de verificación. Cada ítem ahora incluye:

\begin{lstlisting}[caption=Estructura de itemsToCheck en CouchDB, label=lst:items-to-check]
{
    "itemsToCheck": [
        {
            "sku": "C05000",
            "value": 8,
            "verification": "precise",
            "calculationType": "cpu-hours",
            "thresholds": {
                "warning": 0.99,
                "error": 0.975
            }
        },
        {
            "sku": "MKS1000",
            "verification": "existence-only"
        }
    ]
}
\end{lstlisting}

Esta estructura permite definir de forma declarativa qué verificación aplicar a cada SKU, facilitando la adición de nuevos recursos sin modificar el código.

\subsubsection{Endpoint de verificación de cantidades diarias}
Se ha creado un nuevo endpoint específico para la verificación de cantidades diarias:

\begin{verbatim}
GET /data/v1/utilization/check-daily-quantity/:contractId
\end{verbatim}

Este endpoint realiza el siguiente proceso interno:
\begin{enumerate}
    \item Calcula los días transcurridos en el período de facturación actual.
    \item Obtiene los datos de uso desde la billing-api mediante \texttt{/\{contractId\}/usage}.
    \item Agrega el consumo por SKU utilizando la función \texttt{aggregateUsageData()}.
    \item Compara los valores actuales contra los esperados para cada SKU configurado.
\end{enumerate}

\subsubsection{Cálculo del valor esperado}
Para los SKUs con verificación \texttt{precise}, el valor esperado se calcula mediante la fórmula:

\begin{equation}
\texttt{expected} = \texttt{value} \times 24 \text{ horas} \times \texttt{días\_transcurridos}
\end{equation}

Donde \texttt{value} es el valor base configurado (por ejemplo, 8 cores) y \texttt{días\_transcurridos} son los días desde el inicio del período de facturación hasta ayer.

\subsubsection{Agregación de datos de uso}
La función \texttt{aggregateUsageData()} procesa la respuesta de la billing-api y agrupa el consumo por SKU:

\begin{lstlisting}[caption=Función de agregación de datos de uso, label=lst:aggregate-usage]
function aggregateUsageData(usageData) {
    const aggregated = {};
    usageData.datacenters.forEach(dc => {
        dc.meters.forEach(meter => {
            const sku = meter.meterId;
            if (!aggregated[sku]) {
                aggregated[sku] = 0;
            }
            aggregated[sku] += parseFloat(meter.quantity.quantity);
        });
    });
    return aggregated;
}
\end{lstlisting}

\subsubsection{Verificación y comparación}
La función \texttt{processItemVerification()} determina el tipo de verificación a aplicar:

\begin{lstlisting}[caption=Lógica de verificación por tipo, label=lst:process-verification]
if (itemConfig.verification === 'existence-only') {
    return this.generateExistenceVerification(sku, actualValue);
}
// Para verificacion precise:
const expected = itemConfig.value * 24 * daysElapsed;
const ratio = (actualValue / expected) * 100;
return this.evaluateWithThresholds(sku, actualValue, expected, ratio);
\end{lstlisting}

\subsubsection{Respuesta del endpoint}
El endpoint devuelve una respuesta estructurada con el resultado de cada verificación:

\begin{lstlisting}[caption=Ejemplo de respuesta del endpoint de cantidades diarias, label=lst:response-daily]
{
  "ok": true,
  "message": "daily-quantities check completed",
  "contractId": "36348704",
  "period": "this month until yesterday (3 days)",
  "daysCalculated": 3,
  "results": [
    {
      "sku": "C05000",
      "status": "ok",
      "actual": 1920.00,
      "expected": 1920.00,
      "ratio": 100,
      "calculationMethod": "precise"
    }
  ],
  "summary": {
    "total": 10,
    "ok": 8,
    "warnings": 1,
    "errors": 1
  }
}
\end{lstlisting}

El campo \texttt{ratio} se expresa como porcentaje, donde 100 indica coincidencia perfecta entre el valor actual y el esperado.

\subsubsection{Configuración de PostMan con Bearer Token estándar}
Se ha actualizado la colección de PostMan para utilizar la pestaña \textit{Authorization} con tipo \texttt{Bearer Token} en lugar de una variable personalizada en los headers. Este cambio:

\begin{itemize}
    \item Aprovecha la interfaz nativa de PostMan para gestión de tokens.
    \item Permite heredar la autenticación a nivel de colección, evitando configurarla en cada petición.
    \item Facilita el uso por parte de otros miembros del equipo que estén familiarizados con PostMan.
\end{itemize}

\subsubsection{Scheduler para verificación de cantidades diarias}
Se ha configurado un nuevo thread en el scheduler de CouchDB específico para la verificación de cantidades:

\begin{lstlisting}[caption=Configuración del scheduler para cantidades diarias, label=lst:scheduler-quantities]
{
    "thread": "utilization-verifier:daily-quantities",
    "enabled": true,
    "mutex": true,
    "timeout": 3600000,
    "action": "utilization-verifier:daily-quantities:request",
    "endAction": "utilization-verifier:daily-quantities:complete"
}
\end{lstlisting}

Este scheduler puede programarse con expresiones cron independientes del verificador de existencia, permitiendo ejecutar las verificaciones de cantidades con una frecuencia diferente si fuera necesario.



\subsection{Evaluación y reunión con el cliente}
Los cambios solicitados por el cliente se han implementado correctamente, simplificando el uso de la API y mejorando la configuración de PostMan. Además al volver a reunirme con el cliente dio el visto bueno a la iteración y hablamos sobre un posible cambio en el futuro de seprarar las credenciales de las configuraciones en dos documentos/formatos distintos, siendo este comentario más orientado a el grupo de trabajo entero que a mi trabajo, ya que esta decisión es propia del equipo de trabajo.

Referido a los tiempos, ocurre algo similar a la iteración anterior, se han finalizado antes de lo previsto, pero las horas totales dedicadas se han aproximado bastante a la estimación inicial. La desviación en la tarea de implementación de la verificación de cantidades se debe principalmente a la complejidad añadida por el sistema de thresholds y la necesidad de realizar pruebas adicionales para validar los diferentes escenarios (valores dentro del umbral de warning, valores que superan el umbral de error, etc.).

En resumen, la variación en los tiempos de esta iteración se debe a que fue necesario emplear tiempo adicional para implementar las mejoras solicitadas por el cliente, como el sistema de thresholds, la estandarización del contrato, el uso de Bearer Token en PostMan y la mejora del sistema de alertas. Esta complejidad añadida implicó realizar pruebas adicionales para validar diferentes escenarios (valores dentro del umbral de warning, valores que superan el umbral de error, etc.). Sin embargo, gracias a que las pruebas y la validación se simplificaron por estas mismas mejoras, el tiempo total invertido no se vio muy afectado, manteniéndose próximo a la estimación inicial y finalizándose antes de lo previsto.
Además de que se decidio desglosar y plasmar también el tiempo de reunión con el cliente, ya que aunque es una tarea que no aporta valor directo al desarrollo, es necesaria para la correcta evolución del proyecto y para garantizar que se cumplen las expectativas del cliente, por lo que se considera relevante incluirla en la comparativa de horas.

\begin{table}[H]
    \centering
    \begin{tabular}{|l|c|c|c|}
        \hline
        \textbf{Tarea} & \textbf{Estimadas (h)} & \textbf{Reales (h)} & \textbf{Desviación}\\
        \hline
        Verificación cantidades & 16h & 18h & +12.5\% \\
        Comprobación contrato completo & 1h & 10' & -83.3\% \\
        Configuraciones & 2h & 6h 10min & +205\% \\
        Pruebas y validación & 4h & 1h & -75\% \\
        Reunión con el cliente & 0h & 1h & $\infty$\% \\
        \hline
        \textbf{Total} & \textbf{24h} & \textbf{26h} & \textbf{+8.3\%} \\
        \hline
    \end{tabular}
    \caption{Comparativa de horas estimadas vs reales en la Iteración 2}
    \label{tab:horas-iteracion2}
\end{table}

% ============================================
% ===== TERCERA ITERACIÓN =====
% ============================================
\newpage
\section{Iteración 3}

\subsection{Planificación}

En esta iteración se cierra el ciclo de verificación de los \textit{items} continuables iniciado en las iteraciones anteriores: conocida la existencia de los SKUs (\textbf{TFG 3.1.1}) y verificadas sus cantidades de uso (\textbf{TFG 3.2.1}), el siguiente paso natural es contrastar el coste económico teórico con el importe real reflejado en la factura de IONOS Cloud (\textbf{TFG 3.3.1}).

\textbf{Comprobación del gasto de ítems continuables:} A partir de las cantidades de uso obtenidas de la \textit{billing-api} y de las tarifas unitarias configuradas en CouchDB, se calcula el importe esperado para cada SKU y se contrasta con el valor real de la factura. Este proceso resulta fundamental para satisfacer los requisitos RF03.01, RF03.02 y RF03.04, cuya ejecución está también vinculada al requisito RF04.02.

\textbf{Obtención y validación de la factura:} Se consume el endpoint de facturas de la API de IONOS Cloud y se valida la respuesta mediante el esquema Joi definido en iteraciones anteriores, asegurando que únicamente se procesan facturas cerradas y definitivas (\texttt{finallyPosted === true}).

\textbf{Sistema de tolerancias:} Se reutiliza y extiende el sistema de umbrales configurables implementado en la iteración anterior, añadiendo parámetros específicos para la validación de importes monetarios y \textit{rebates}.

Como en las iteraciones precedentes, se repiten las tareas transversales: verificación de completitud del contrato (\textbf{TFG 3.3.2}), actualización de credenciales y configuración en CouchDB (\textbf{TFG 3.3.3}), pruebas de funcionamiento (\textbf{TFG 3.3.4}) y configuración de alertas y rutinas (\textbf{TFG 3.3.5}).

\subsection{Análisis}

\subsubsection{Obtención de la factura real}

La \textit{billing-api} expone dos endpoints para el acceso a los datos de facturación:

\begin{enumerate}
    \item \textbf{Listado de facturas:} \texttt{GET /\{contractId\}/invoices} --- Devuelve la lista básica de facturas disponibles con su identificador y período de facturación.
    \item \textbf{Detalle de factura:} \texttt{GET /\{contractId\}/invoices/\{invoiceId\}} --- Devuelve la factura completa con el desglose por \textit{datacenter} y SKU.
\end{enumerate}

Los campos más relevantes de la respuesta son:

\begin{itemize}
    \item \textbf{\texttt{metadata.postingPeriod}:} Período de facturación en formato \texttt{YYYY-MM}.
    \item \textbf{\texttt{metadata.finallyPosted}:} Booleano que indica si la factura está cerrada y definitiva.
    \item \textbf{\texttt{total.quantity} / \texttt{total.unit}:} Importe total de la factura.
    \item \textbf{\texttt{datacenters[].meters[]}:} Desglose por SKU, con los campos \texttt{meterId}, \texttt{quantity}, \texttt{rate} y \texttt{amount}.
    \item \textbf{\texttt{datacenters[].rebate.amount.quantity}:} Descuento aplicado por \textit{datacenter}.
\end{itemize}

El verificador trabaja siempre sobre el \textbf{mes anterior cerrado}: se calcula el período del mes anterior y se busca en el listado la factura cuyo campo \texttt{date} coincida. El campo \texttt{finallyPosted} no está disponible en el endpoint de listado, por lo que se comprueba tras obtener el detalle completo de la factura mediante una llamada separada a \texttt{getInvoiceById()}. Si la factura aún no está disponible o \texttt{finallyPosted} es \texttt{false}, el proceso finaliza sin generar alertas, ya que podría seguir modificándose durante los primeros días del mes.

\subsubsection{Cálculo del coste teórico}

La fórmula de cálculo es directa. Para cada SKU se multiplica la cantidad de uso reportada por la tarifa unitaria almacenada en CouchDB bajo la clave \texttt{skuPrices}:

\begin{equation}
    \texttt{expectedAmount} = \texttt{expectedQuantity} \times \texttt{expectedRate}
\end{equation}

La tarifa \texttt{expectedRate} es independiente de la que figura en la factura. Esto es una decisión de diseño deliberada: si los precios de la factura contuviesen algún error, usar esa misma fuente para calcular el valor de referencia invalidaría cualquier detección de anomalía. Al mantenerla separada en CouchDB, el sistema conserva una fuente de verdad propia.

Las unidades se utilizan tal cual, sin conversiones adicionales: si la tarifa está expresada en EUR/30Days, la cantidad de uso vendrá igualmente en esas unidades.

\subsubsection{Sistema de tolerancias y estados}

Se extiende el sistema de umbrales de la iteración anterior con parámetros específicos para la validación de cantidades e importes. Los valores por defecto, configurables por SKU en CouchDB bajo \texttt{config.invoiceValidation.validationTolerances}, son los siguientes:

\begin{table}[H]
    \centering
    \begin{tabular}{|p{0.40\textwidth}|p{0.18\textwidth}|p{0.35\textwidth}|}
        \hline
        \textbf{Parámetro} & \textbf{Valor por defecto} & \textbf{Descripción} \\
        \hline
        \hline
        \texttt{DEFAULT\_QUANTITY\_TOLERANCE} & 1\% (0,01) & Tolerancia general de cantidad \\
        \hline
        \texttt{DEFAULT\_ERROR\_THRESHOLD} & 95\% (0,95) & Ratio mínimo para no generar error en cantidades \\
        \hline
    \end{tabular}
    \caption{Parámetros de tolerancia para validación de facturas}
    \label{tab:tolerancias-factura}
\end{table}

La lógica de estados difiere según el tipo de comparación:

\begin{itemize}
    \item \textbf{Validación de cantidades} (\texttt{calculateQuantityComparison}): tres estados posibles según el ratio respecto al valor esperado.
    \begin{itemize}
        \item \textbf{\texttt{ok}:} ratio $\geq$ 0,99 --- coincidencia prácticamente perfecta.
        \item \textbf{\texttt{warning}:} 0,95 $\leq$ ratio $<$ 0,99 --- desviación menor, requiere atención.
        \item \textbf{\texttt{error}:} ratio $<$ 0,95 --- discrepancia crítica, requiere acción inmediata.
    \end{itemize}
    \item \textbf{Validación de importes monetarios} (\texttt{validateAmounts}): únicamente dos estados: \textbf{\texttt{ok}} si el importe está dentro del umbral configurado, \textbf{\texttt{error}} en caso contrario. No existe estado intermedio \texttt{warning} para importes.
\end{itemize}

\subsection{Desarrollo y pruebas}

\subsubsection{Obtención y validación de la factura}

El primer paso del proceso es localizar la factura correspondiente al mes anterior cerrado. Para ello se utiliza el helper \texttt{InvoiceHelper}, que encapsula las llamadas a la \textit{billing-api}. El endpoint de listado devuelve únicamente metadatos básicos — entre ellos el campo \texttt{date} con el período — pero \textbf{no incluye \texttt{finallyPosted}}. Por ello, la búsqueda inicial filtra solo por \texttt{date}; la comprobación de \texttt{finallyPosted} se realiza tras obtener el detalle completo con \texttt{getInvoiceById()}:

\begin{lstlisting}[caption=Obtención de la factura del mes anterior cerrado, label=lst:get-invoice]
const targetPeriod = now.clone().subtract(1, 'month')
                                .startOf('month')
                                .format('YYYY-MM');

// NOTE: The list endpoint does NOT include finallyPosted,
// so we must fetch details to check it.
const invoices       = await invoiceHelper.listInvoices(contractId, token);
const matchingInvoice = invoices.find((inv) => inv.date === targetPeriod);
\end{lstlisting}

Si no se localiza ninguna factura para ese período, o si tras obtener el detalle \texttt{finallyPosted} es \texttt{false}, el verificador termina sin generar alertas. A continuación se valida la respuesta mediante \texttt{invoiceSchema} antes de continuar con el procesamiento. De este modo se garantiza que los campos esperados están presentes y con el formato correcto, evitando errores de acceso en etapas posteriores.

\subsubsection{Cálculo del coste teórico por SKU}

Con la factura validada se recorre el array de \textit{datacenters} y sus \textit{meters}. Para cada SKU se calcula el importe esperado y se determina si la diferencia absoluta respecto al importe real queda dentro de la tolerancia configurada:

\begin{lstlisting}[caption=Cálculo del importe esperado y validación de la desviación, label=lst:expected-amount]
const expectedRate    = skuPrices[meter.meterId];
const expectedAmount  = expectedQuantity * expectedRate;
const invoiceAmount   = parseFloat(meter.amount.quantity);

const difference      = invoiceAmount - expectedAmount;
const toleranceAmount = Math.abs(expectedAmount) * DEFAULT_QUANTITY_TOLERANCE;
const isValid         = Math.abs(difference) <= Math.max(toleranceAmount, REBATE_TOLERANCE);
\end{lstlisting}


\subsubsection{Validación de \textit{rebates}}

Los descuentos por \textit{datacenter} se validan de forma independiente al resto de la jerarquía, usando una tolerancia absoluta de 0,01. Esta tolerancia es más estricta en términos absolutos que la tolerancia porcentual de los SKUs, dado que los \textit{rebates} son importes fijos y pequeñas diferencias de redondeo son esperables en su cálculo:

\begin{lstlisting}[caption=Validación del rebate por datacenter, label=lst:rebate]
const rebateDiff   = Math.abs(actualRebate - expectedRebate);
const rebateStatus = rebateDiff < REBATE_TOLERANCE
    ? 'ok'
    : 'error';
\end{lstlisting}

\subsubsection{Endpoint de validación y ejecución programada}

Se ha añadido un endpoint de prueba manual que permite lanzar la validación sobre un contrato concreto sin esperar a la ejecución programada:

\begin{verbatim}
POST /data/v1/invoice/manual-test/:contractId
\end{verbatim}

Este endpoint devuelve una respuesta estructurada con el resultado por SKU, por grupo de producto y el estado global del contrato, lo que facilita la depuración y la revisión puntual de facturas.

Para la ejecución automática se ha añadido un nuevo \textit{thread} al \textit{scheduler} de CouchDB. El verificador corre los días 3, 4 y 5 de cada mes a las 14:00 (Europe/Berlin), dando margen para que la factura del mes anterior esté cerrada y disponible. La opción \texttt{mutex: true} garantiza que no se ejecutan dos instancias en paralelo si una tarda más de lo esperado:

\begin{lstlisting}[caption=Configuración del scheduler para validación de facturas, label=lst:scheduler-invoice]
{
    "thread":    "invoice-verifier:monthly",
    "enabled":   true,
    "mutex":     true,
    "timeout":   3600000,
    "action":    "invoice-verifier:monthly:request",
    "endAction": "invoice-verifier:monthly:complete"
}
\end{lstlisting}

Esta configuración satisface los requisitos RF04.01 (ejecución programada del sistema de validación) y RF03.04 (programación de \textit{jobs} cron), ya que el ciclo mensual encaja con la periodicidad de las facturas de IONOS Cloud.

\subsubsection{Comprobación de completitud del contrato (TFG 3.3.2)}

Como en las iteraciones anteriores, antes de comenzar con el desarrollo de la funcionalidad principal se verifica que el contrato de pruebas sigue conteniendo todos los SKUs necesarios (\textbf{TFG 3.3.2}). En este caso no fue necesario añadir ningún elemento nuevo, ya que los SKUs monitorizados no habían variado respecto a la iteración anterior.

\subsubsection{Actualización de credenciales y configuración (TFG 3.3.3)}

Se actualizó el documento de configuración de CouchDB para incluir los nuevos parámetros de tolerancia de facturación (\textbf{TFG 3.3.3}): \texttt{validationTolerances}, \texttt{skuPrices} y la configuración del nuevo \textit{thread} del \textit{scheduler}. La estructura de credenciales no requirió cambios adicionales, ya que el helper \texttt{InvoiceHelper} reutiliza los mismos tokens de acceso definidos en iteraciones previas.

\subsubsection{Pruebas manuales con PostMan (TFG 3.3.4)}

Se añadieron nuevas colecciones en PostMan para la validación de facturas (\textbf{TFG 3.3.4}), verificando los siguientes escenarios:

\begin{itemize}
    \item Respuesta 200 cuando la factura del mes anterior está disponible y finalizada.
    \item Estado global \texttt{ok} cuando todos los importes coinciden dentro del umbral configurado.
    \item Paso a \texttt{error} de un SKU concreto al introducir una tarifa incorrecta en \texttt{skuPrices} de CouchDB que supera la tolerancia configurada (la validación de importes solo produce \texttt{ok} o \texttt{error}, sin estado intermedio).
    \item Finalización sin alertas cuando \texttt{finallyPosted === false}, comprobando que el sistema espera a que la factura esté cerrada.
\end{itemize}

\subsubsection{Configuración de alertas (TFG 3.3.5)}

Se adaptó el sistema de alertas para incluir el nuevo verificador de facturas (\textbf{TFG 3.3.5}). Las alertas generadas siguen la misma estructura que las de iteraciones anteriores: se notifica el estado global del contrato y, en caso de \texttt{warning} o \texttt{error}, se detallan los SKUs y \textit{datacenters} afectados junto con el ratio calculado. Esto permite al equipo de operaciones identificar rápidamente si la discrepancia corresponde a un SKU concreto, a un grupo de producto o a un \textit{rebate} incorrecto.

\subsection{Evaluación}

En esta iteración se ha implementado la validación económica completa de los \textit{items} continuables, cerrando el ciclo de verificación que comenzó en la iteración 1 (detección de existencia de SKUs) y se amplió en la iteración 2 (verificación de cantidades de uso). El sistema puede ahora comparar el gasto teórico calculado con el importe real facturado, desglosándolo de forma jerárquica por SKU, grupo de producto y \textit{datacenter}, y generando alertas diferenciadas por nivel de criticidad.

Con esta iteración quedan satisfechos los requisitos \textbf{RF03.01} (cálculo del coste esperado por SKU), \textbf{RF03.02} (comparación con el importe real de la factura), \textbf{RF03.04} (programación de la ejecución mediante cron) y \textbf{RF04.02} (contraste entre el cálculo teórico y los datos reales de la API).

En esta reunión no fue necesaria la reunión con el cliente, ya que los cambios introducidos en esta iteración fueron pequeños y con la comunicación de los cambios y el modo en los que se hicieros via mensaje fue suficiente y no se requirió reunión. 

Referido a tiempos, se realizo según lo previsto aproximadamente, invirtiendo las horas necesarias, necesitanto menos tiempo para comprobar que el contrato estaba compledo debido a que no se introdujeron ningún cambio.

\begin{table}[H]
    \centering
    \begin{tabular}{|l|c|c|c|}
        \hline
        \textbf{Tarea} & \textbf{Estimadas (h)} & \textbf{Reales (h)} & \textbf{Desviación} \\
        \hline
        Comprobación gasto \textit{items} continuables & 16h & 16h & 0\% \\
        Comprobación contrato completo & 1h & 10' & -83\% \\
        Configuración de credenciales, alertas y rutinas & 2h & 2h y 30' & +25\% \\
        Testeo de funcionamiento & 4h & 4h & 0\% \\
        \hline
        \textbf{Total} & \textbf{23h} & \textbf{22h y 40'} & \textbf{3\%} \\
        \hline
    \end{tabular}
    \caption{Comparativa de horas estimadas vs reales en la Iteración 3}
    \label{tab:horas-iteracion3}
\end{table}

% ============================================
% ===== CUARTA ITERACIÓN =====
% ============================================
\newpage
\section{Iteración 4}

\subsection{Planificación}

Las tres iteraciones anteriores han completado el ciclo de verificación de los \textit{items} continuables: se detecta su presencia en el contrato, se contrasta su consumo diario y se valida su coste en la factura mensual. En esta cuarta iteración el foco cambia: se abordan por primera vez los \textit{items} no continuables, representados por los tokens de IA de IONOS Cloud (\textbf{TFG 3.4.1}).

A diferencia de los recursos de infraestructura (cores de CPU, RAM o almacenamiento), el consumo de tokens de IA solo existe si el sistema realiza llamadas activas al endpoint de inferencia. Sin generación controlada, la API de facturación simplemente no reflejaría consumo para los SKUs \texttt{AIL370I1100} (tokens de entrada) y \texttt{AIL370O1100} (tokens de salida), haciendo imposible cualquier verificación.

\textbf{Módulo generador de tokens:} Se implementará un módulo \texttt{generator} que, activado diariamente por el \textit{scheduler}, realice llamadas repetidas al endpoint de IA de IONOS Cloud hasta acumular un número mínimo de tokens configurable (\texttt{targetTokens}). El diseño sigue el mismo patrón basado en eventos ya utilizado en módulos anteriores.

\textbf{Control mediante configuración:} Todos los parámetros relevantes (modelo, \texttt{targetTokens}, \textit{prompt} y URL del endpoint) se centralizan en CouchDB, de modo que pueden ajustarse sin modificar el código.

\subsection{Análisis}

\subsubsection{El problema de los \textit{items} no continuables}

Los \textit{items} continuables ---como los cores de CPU o los GB de RAM--- generan consumo de forma automática mientras los recursos están activos. Los tokens de IA, en cambio, solo aparecen en la factura si se realizan llamadas explícitas al modelo. Esto plantea un problema de verificabilidad: sin consumo no hay datos que verificar, y sin generación controlada no es posible predecir cuántos tokens deberían aparecer.

La solución adoptada es crear un generador que se ejecute a hora fija cada día y produzca un volumen de tokens conocido de antemano. Una vez que el consumo es predecible, el módulo de verificación de utilización puede comprobar los SKUs de IA con verificación \texttt{existence-only}: basta con confirmar que hay consumo mayor que cero en el período.

\subsubsection{Separación de tokens de entrada y salida}

La API de IONOS Cloud distingue dos SKUs para el consumo de IA:

\begin{itemize}
    \item \textbf{\texttt{AIL370I1100}} (tokens de entrada / \textit{input}): corresponde a los tokens del \textit{prompt} enviado al modelo en cada llamada.
    \item \textbf{\texttt{AIL370O1100}} (tokens de salida / \textit{output}): corresponde a los tokens generados por el modelo en la respuesta.
\end{itemize}

Con un \texttt{targetTokens} de 100.000, la distribución habitual es aproximadamente 10.000 de entrada y 90.000 de salida, es decir, 0,01M y 0,09M respectivamente en las unidades de facturación. El \textit{prompt} largo es intencionado: cuanto más extenso sea, más tokens de salida genera el modelo por llamada, reduciendo el número de iteraciones necesarias para alcanzar el objetivo.

\subsubsection{Compatibilidad con el SDK de OpenAI}

El endpoint de inferencia de IONOS Cloud implementa la misma interfaz que la API de OpenAI, lo que permite utilizar el SDK oficial de Node.js (\texttt{openai}) apuntando a la URL del servicio de IONOS. La única diferencia respecto a un uso convencional del SDK es que \texttt{baseURL} apunta al endpoint de IONOS en lugar del de OpenAI, por lo que no fue necesario aprender ni integrar ningún cliente adicional.

El uso de \texttt{temperature: 0} garantiza respuestas deterministas: con el mismo \textit{prompt} y el mismo modelo, el número de tokens generados es reproducible entre ejecuciones, lo que facilita la predicción del consumo mensual acumulado.

\subsection{Desarrollo y pruebas}

\subsubsection{Configuración del \textit{scheduler} en CouchDB}

El generador se activa cada día a las 08:00 mediante un nuevo \textit{thread} en el \textit{scheduler}. La opción \texttt{mutex: true} impide que dos ejecuciones se solapen si la generación tarda más de lo habitual:

\begin{lstlisting}[caption=Configuración del scheduler para el generador de tokens, label=lst:scheduler-generator]
{
    "thread":    "generator",
    "enabled":   true,
    "mutex":     true,
    "timeout":   3600000,
    "action":    "generator:request",
    "endAction": "generator:complete",
    "schedule":  [["0 8 * * *", {}]]
}
\end{lstlisting}

El campo \texttt{action} debe coincidir exactamente con el nombre del evento registrado en el módulo, ya que es la cadena que el \textit{mediator} usa para enrutar el disparo.

\subsubsection{Módulo Generator: registro del evento}

El módulo \texttt{generator} sigue el mismo patrón que los módulos anteriores: en su método \texttt{init()} registra un \textit{listener} sobre el evento correspondiente y delega toda la lógica en el \textit{helper} especializado:

\begin{lstlisting}[caption=Registro del evento en el módulo Generator, label=lst:generator-event]
// src/modules/generator/index.js:22
this.mediator.on(Actions.GENERATOR_REQUEST, async () => {
    const result = await this.inferenceHelper.generateTokens();
    this.mediator.emit(Actions.GENERATOR_COMPLETE); // libera el mutex
});
\end{lstlisting}

La constante \texttt{Actions.GENERATOR\_REQUEST} vale \texttt{'generator:request'}, que coincide exactamente con el \texttt{action} configurado en el \textit{scheduler}. La emisión de \texttt{GENERATOR\_COMPLETE} al finalizar, independientemente del resultado, es lo que libera el \textit{mutex} y permite que la siguiente ejecución programada pueda arrancar.

\subsubsection{InferenceHelper: carga de configuración}

Antes de iniciar la generación, \texttt{InferenceHelper} lee sus parámetros desde CouchDB. Si alguno no está definido, aplica un valor por defecto:

\begin{lstlisting}[caption=Carga de parámetros de configuración en InferenceHelper, label=lst:inference-config]
// src/helpers/inference.js:35-39
const { baseUrl, apiKey, model, targetTokens, prompt }
    = this.getConfig().config || {};

this.targetTokens = targetTokens || 100000;
this.prompt       = prompt       || DEFAULT_PROMPT;
\end{lstlisting}

El \texttt{DEFAULT\_PROMPT} es un texto extenso sobre computación en la nube que actúa como \textit{fallback}. Su longitud es deliberada: un \textit{prompt} largo produce más tokens de salida por llamada, reduciendo el número de iteraciones necesarias para alcanzar \texttt{targetTokens}.

\subsubsection{InferenceHelper: bucle de generación}

El núcleo del módulo es un bucle que acumula tokens hasta superar el umbral configurado. En cada iteración se realiza una única llamada al modelo y se suman los tokens consumidos, separando los de entrada de los de salida:

\begin{lstlisting}[caption=Bucle de generación de tokens hasta alcanzar el objetivo, label=lst:inference-loop]
// src/helpers/inference.js:59
while (inputTokens + outputTokens < this.targetTokens) {
    const response = await this.client.chat.completions.create({
        model:       this.model,
        temperature: 0,
        messages:    [{ role: 'user', content: prompt }]
    });
    inputTokens  += response.usage.prompt_tokens     || 0;
    outputTokens += response.usage.completion_tokens || 0;
    calls++;
}
return { inputTokens, outputTokens,
         totalTokens: inputTokens + outputTokens, calls };
\end{lstlisting}

Con \texttt{targetTokens = 100.000} y un \textit{prompt} suficientemente largo, en la práctica suelen necesitarse una o dos llamadas para alcanzar el objetivo. Un resultado típico sería \texttt{\{inputTokens: 10000, outputTokens: 90000, totalTokens: 100000, calls: 1\}}.

\subsubsection{Actualización del contrato y configuración de verificación}

Los SKUs \texttt{AIL370I1100} y \texttt{AIL370O1100} son nuevos en el contrato, por lo que fue necesario añadirlos. A diferencia de los SKUs continuables, se configuran con verificación \texttt{existence-only}: no se compara contra un valor esperado fijo, sino que basta con confirmar que hay consumo mayor que cero, ya que la cantidad exacta puede variar ligeramente según el número de días transcurridos en el período:

\begin{lstlisting}[caption=Configuración de los SKUs de IA en itemsToCheck, label=lst:ai-sku-config]
{ "sku": "AIL370I1100", "verification": "existence-only" },
{ "sku": "AIL370O1100", "verification": "existence-only" }
\end{lstlisting}

\subsection{Evaluación}

Con esta iteración el sistema es capaz de generar consumo de IA de forma autónoma y verificar que dicho consumo aparece en la API de facturación. El módulo generador cierra el ciclo para los \textit{items} no continuables: primero produce el consumo y, en la verificación diaria de cantidades implementada en la iteración 2, se comprueba que los SKUs de IA están presentes.

No se incluyó una subsección de pruebas con PostMan independiente porque el módulo generador no expone un endpoint REST propio: su ejecución está diseñada para ser disparada exclusivamente por el \textit{scheduler}. Para la validación manual se empleó directamente el endpoint de IA de IONOS, comprobando que las llamadas devolvían el campo \texttt{usage} con los contadores esperados antes de integrar la lógica en el bucle.

Del mismo modo,aunque no fue necesaria una reunión con el cliente, porque los cambios se comunicaron de forma asíncrona y no requirieron validación, se hizo una reunión de seguimiento para tratar los posibles porblemas y decisiones que no se habian especificado, el cliente comunico el problema de usar un gasto de tockes lo suficientemente grande para que se plasmará en la factura pero no un tamaño que fuese demasiado grande, por lo que se llego a la conclusión de usar 100.000 tockens.

La implementación del \textit{helper} resultó más directa de lo previsto gracias a la compatibilidad del SDK de OpenAI con el endpoint de IONOS. La mayor inversión de tiempo se concentró en las pruebas de integración, ya que el consumo real de tokens tarda unas horas en reflejarse en la API de facturación y para verificar que funcionaba correctamente fue necesario esperar a que el consumo generado por el módulo se viera reflejado en la API de facturación, lo que alargó el proceso de validación. Sin embargo, una vez superada esta fase, la integración con el sistema de verificación existente fue fluida, ya que los SKUs de IA se configuraron con verificación \texttt{existence-only}, lo que simplificó la lógica de validación. Además los tiempos muertos de espera se usararon para redactar.

\begin{table}[H]
    \centering
    \begin{tabular}{|l|c|c|c|}
        \hline
        \textbf{Tarea} & \textbf{Estimadas (h)} & \textbf{Reales (h)} & \textbf{Desviación} \\
        \hline
        Implementación módulo generador     & 24h & 20h  & -16.7\% \\
        Comprobación contrato completo      & 1h  & 30'  & -50\%   \\
        Configuración de credenciales y rutinas & 4h & 5h & +25\%  \\
        Testeo de funcionamiento            & 8h  & 11h y 30' & +45\%   \\
        Reunión con el cliente                   & 0h  & 30'   & $\infty$ \%   \\
        \hline
        \textbf{Total} & \textbf{37h} & \textbf{37h 30'} & \textbf{+1.3\%} \\
        \hline
    \end{tabular}
    \caption{Comparativa de horas estimadas vs reales en la Iteración 4}
    \label{tab:horas-iteracion4}
\end{table}

% ============================================
% ===== QUINTA ITERACIÓN =====
% ============================================
\newpage
\section{Iteración 5}

\subsection{Planificación}

Esta iteración cierra el ciclo completo del proyecto: tras haber generado consumo de tokens de IA en la iteración anterior y haber verificado su existencia en la API de utilización, queda pendiente comprobar que ese consumo se traduce correctamente en la \textbf{factura mensual} de IONOS Cloud. El objetivo es validar el coste de los SKUs \texttt{AIL370I1100} (tokens de entrada) y \texttt{AIL370O1100} (tokens de salida) frente al valor teórico esperado (\textbf{TFG 3.5.1}).

\textbf{Reutilización de la infraestructura de la Iteración 3:} no es necesario implementar un nuevo verificador, sino extender el ya existente. Se reaprovechan \texttt{InvoiceHelper}, el esquema Joi \texttt{invoiceSchema} y el \textit{thread} \texttt{invoice-verifier:monthly} del \textit{scheduler}, lo que reduce el esfuerzo de la iteración a la incorporación de un nuevo tipo de verificación y a su configuración en CouchDB.

\textbf{Diferencia clave respecto a la Iteración 3:} los tokens son \textit{items} no continuables, por lo que requieren un sistema de tolerancias distinto del aplicado a los \textit{items} continuables. En particular, los tokens de salida presentan una variabilidad inherente al modelo que obliga a admitir un margen mayor que el de cantidades continuas o importes monetarios.

\subsection{Análisis}

\subsubsection{Tarifas y cálculo del valor esperado}

El primer paso es conocer la tarifa real aplicada por IONOS Cloud a cada SKU de IA. Esta información se obtiene consultando la \textit{billing-api} con el filtro \texttt{meterDesc} sobre los identificadores \texttt{AIL370I1100} y \texttt{AIL370O1100}, y se almacena en CouchDB bajo \texttt{skuPrices} para mantener una fuente de verdad propia, independiente de la factura, tal y como se decidió en la iteración 3.

A partir del consumo diario fijado en la iteración 4 (\texttt{targetTokens = 100.000}, distribuidos típicamente como un 10\% de \textit{input} y un 90\% de \textit{output}), la cantidad esperada al final del mes se calcula mediante la siguiente expresión:

\begin{equation}
    \texttt{expectedQuantity} = \frac{\texttt{targetTokens diarios} \times \texttt{daysPassed}}{1.000.000}
\end{equation}

La división por un millón se debe a que la facturación de IONOS Cloud expresa estos SKUs en unidades de millones de tokens. Para un mes de 30 días, los valores de referencia resultan ser aproximadamente \texttt{0{,}3M} para \texttt{AIL370I1100} y \texttt{2{,}7M} para \texttt{AIL370O1100}.

\subsubsection{Sistema de tolerancias para \textit{items} no continuables}

El planteamiento de tolerancias de los \textit{items} continuables (cantidades constantes y predecibles dentro del 1\%) no se aplica directamente aquí. Los tokens son \textit{items} no continuables y por su naturaleza requieren un margen distinto:

\begin{itemize}
    \item \textbf{\texttt{AIL370I1100} (tokens de entrada):} tolerancia mínima, próxima a 0. El \textit{prompt} es fijo y el SDK siempre tokeniza de la misma manera, por lo que el número de tokens de entrada es prácticamente determinista.
    \item \textbf{\texttt{AIL370O1100} (tokens de salida):} tolerancia mayor, en torno al 5--10\%. Aunque la generación se ejecuta con \texttt{temperature: 0}, el modelo puede producir respuestas ligeramente distintas en versiones sucesivas o ante cambios menores en el \textit{backend}, por lo que conviene admitir cierto margen sin que ello dispare una alerta.
\end{itemize}

Para soportar esta lógica diferenciada se reutiliza el tipo de verificación \texttt{non-consumable}, ya disponible en el módulo \texttt{utilization} pero hasta el momento sin uso real, ya que en la iteración 4 los SKUs de IA se configuraron con \texttt{existence-only}. Cada SKU declara dos parámetros: \texttt{value} (cantidad diaria esperada) y \texttt{tolerance} (margen aceptable expresado como fracción).

\subsubsection{Dispatcher de tipos de verificación}

La extensión del sistema se apoya en el patrón \textit{Strategy} ya presente en \texttt{processItemVerification()}, que enruta cada SKU a la función adecuada según su campo \texttt{verification}:

\begin{lstlisting}[caption=Dispatcher de tipos de verificación en el módulo utilization, label=lst:dispatcher-verif]
async processItemVerification(itemConfig, usageData, daysPassed) {
    const sku         = itemConfig.sku;
    const actualValue = usageData[sku] || 0;

    if (itemConfig.verification === 'existence-only') {
        return this.generateExistenceVerification(sku, actualValue);
    }
    if (itemConfig.verification === 'non-consumable') {
        return this.generateNonConsumableVerification(
            itemConfig, actualValue, daysPassed);
    }
    return this.generatePreciseVerification(itemConfig, actualValue, daysPassed);
}
\end{lstlisting}

Añadir un nuevo tipo de verificación se reduce a incorporar una rama \texttt{if} y la función correspondiente, sin tocar el resto del módulo. Este diseño contribuye al cumplimiento del requisito \textbf{RNF02} (mantenibilidad), ya que la lógica específica de cada estrategia queda aislada en su propia función.


\subsection{Desarrollo y pruebas}

\subsubsection{Configuración en CouchDB}

La configuración de los nuevos SKUs en CouchDB recoge tanto la tarifa unitaria como los parámetros del tipo de verificación \texttt{non-consumable}:

\begin{lstlisting}[caption=Configuración de los SKUs de IA para la validación de factura, label=lst:ai-invoice-config]
{
    "itemsToCheck": [
        { "sku": "AIL370I1100", "verification": "non-consumable",
          "value": 10000, "tolerance": 0 },
        { "sku": "AIL370O1100", "verification": "non-consumable",
          "value": 90000, "tolerance": 0.1 }
    ],
    "skuPrices": {
        "AIL370I1100": 0.50,
        "AIL370O1100": 1.50
    }
}
\end{lstlisting}

El SKU de \textit{input} declara \texttt{tolerance: 0} porque, dado el determinismo del proceso (\texttt{temperature: 0} y \textit{prompt} fijo), no se justifica margen alguno. El SKU de \textit{output}, en cambio, admite un \texttt{tolerance: 0.1} (\textit{$\pm$10\%}) que absorbe la pequeña variabilidad del modelo entre versiones sin generar falsos positivos.

\subsubsection{Función \texttt{generateNonConsumableVerification}}

La función núcleo de la iteración recibe la configuración del \textit{item}, el valor real observado y los días transcurridos en el período. A partir de ahí calcula la desviación y selecciona el estado correspondiente:

\begin{lstlisting}[caption=Verificación de items no continuables, label=lst:non-consumable]
generateNonConsumableVerification(itemConfig, actualValue, daysPassed) {
    const expected  = itemConfig.value * daysPassed;
    const tolerance = itemConfig.tolerance ?? 0;

    if (expected === 0 && actualValue === 0) {
        return { sku: itemConfig.sku, status: 'ok', /* ... */ };
    }

    const ratio     = expected > 0 ? actualValue / expected : Infinity;
    const deviation = Math.abs(ratio - 1);

    let status = 'ok';
    if (deviation > tolerance * Utilization.NON_CONSUMABLE_ERROR_MULTIPLIER) {
        status = 'error';
    } else if (deviation > tolerance) {
        status = 'warning';
    }

    return { sku: itemConfig.sku, status, expected, actual: actualValue,
             ratio: Math.round(ratio * 100), tolerance, /* ... */ };
}
\end{lstlisting}

Tres detalles merecen comentario:

\begin{itemize}
    \item El caso borde \texttt{expected === 0 \&\& actualValue === 0} se contempla explícitamente para evitar la división por cero al inicio del mes, cuando \texttt{daysPassed} aún puede ser pequeño y no se ha generado consumo medible.
    \item La fórmula \texttt{deviation = |ratio - 1|} permite detectar desviaciones \textbf{bidireccionales}: tanto un consumo facturado superior al esperado como uno inferior dispararán una alerta, lo cual es deseable porque ambas situaciones indican una posible incidencia.
    \item La constante \texttt{NON\_CONSUMABLE\_ERROR\_MULTIPLIER}, fijada en 2, separa de manera explícita los umbrales de \texttt{warning} y \texttt{error}, lo que permite ajustar la sensibilidad sin tener que tocar la lógica de la función.
\end{itemize}

\subsubsection{Integración con la validación de facturas}

El módulo \texttt{invoice-validation} extrae los \textit{meters} de IA de la respuesta de la \textit{billing-api} y, a través del \textit{dispatcher} comentado en el análisis, los enruta a \texttt{generateNonConsumableVerification()}. El resto del flujo es compartido con la iteración 3: se valida la factura con \texttt{invoiceSchema}, se aplica la tolerancia configurada y se acumulan los resultados en la respuesta jerárquica. Los SKUs que no estén declarados en \texttt{itemsToCheck} pero sí aparezcan en la factura se agregan en una lista \texttt{unmappedSkus}, lo que evita que un SKU nuevo pase desapercibido.

\subsubsection{Comprobación del contrato y configuración (TFG 3.5.2 y 3.5.3)}

Los SKUs \texttt{AIL370I1100} y \texttt{AIL370O1100} ya estaban incluidos en el contrato desde la iteración 4, por lo que no fue necesario añadirlos (\textbf{TFG 3.5.2}). La actualización de configuración en CouchDB se limitó a las nuevas entradas en \texttt{itemsToCheck} y a los precios en \texttt{skuPrices} (\textbf{TFG 3.5.3}). El \textit{thread} \texttt{invoice-verifier:monthly} del \textit{scheduler} se reutiliza tal cual, sin modificaciones, ya que la periodicidad mensual sigue siendo la adecuada.

\subsubsection{Pruebas manuales con PostMan (TFG 3.5.4)}

Las pruebas se ejecutaron una vez cerrado el período mensual, lanzando la validación a través del endpoint manual ya disponible:

\begin{verbatim}
POST /data/v1/invoice/manual-test/:contractId
\end{verbatim}

Se cubrieron cuatro escenarios (\textbf{TFG 3.5.4}):

\begin{itemize}
    \item \textbf{Caso \texttt{ok}:} la cantidad facturada queda dentro de la tolerancia configurada. Es el comportamiento esperado en condiciones normales.
    \item \textbf{Caso \texttt{warning}:} desviación entre 1$\times$ y 2$\times$ la tolerancia. Se forzó modificando ligeramente el campo \texttt{value} de \texttt{AIL370O1100} en CouchDB para situar el ratio en torno al 109\% del valor esperado.
    \item \textbf{Caso \texttt{error} por desviación:} se modificó \texttt{value} para que el cociente quedase muy por encima del doble de la tolerancia, comprobando la transición de \texttt{warning} a \texttt{error}.
    \item \textbf{Caso \texttt{error} por tarifa incorrecta:} se introdujo un valor erróneo en \texttt{skuPrices}, provocando que el cálculo del importe esperado divergiese del importe facturado más allá de la tolerancia.
\end{itemize}

A modo de ejemplo, la respuesta del endpoint para un caso \texttt{warning} en \texttt{AIL370O1100} adopta la siguiente forma:

\begin{lstlisting}[caption=Ejemplo de respuesta del verificador para un SKU no continuable, label=lst:resp-non-consumable]
{
  "sku": "AIL370O1100",
  "status": "warning",
  "expected": 2700000,
  "actual": 2950000,
  "ratio": 109,
  "calculationMethod": "non-consumable",
  "tolerance": 0.1,
  "description": "Non-consumable item - tolerance: +/-10%"
}
\end{lstlisting}

El campo \texttt{calculationMethod} se incluye en la respuesta precisamente para que el equipo de operaciones pueda distinguir, de un vistazo, qué tipo de verificación ha generado cada resultado.

\subsubsection{Configuración de alertas (TFG 3.5.5)}

El sistema de alertas se adapta para incorporar los SKUs de IA dentro del mismo \textit{payload} de Google Chat ya empleado en iteraciones anteriores (\textbf{TFG 3.5.5}). La diferencia principal es que, junto al \texttt{status} y al SKU, se muestra explícitamente el campo \texttt{ratio} (en porcentaje) y la \texttt{tolerance} aplicada. Esta información permite al equipo distinguir rápidamente si una alerta corresponde a una desviación marginal de \texttt{output} o a un problema más serio, como una tarifa mal configurada.

\subsection{Evaluación}

Con esta iteración el sistema queda completo: valida tanto los \textit{items} continuables como los no continuables, en el doble plano del uso diario (módulo \texttt{utilization}) y de la facturación mensual (módulo \texttt{invoice-validation}). Quedan así cubiertos los requisitos \textbf{RF03.01} (cálculo del coste esperado), \textbf{RF03.02} (comparación con el importe real), \textbf{RF03.03} (tolerancia configurable, especialmente relevante en esta iteración por la naturaleza de los tokens), \textbf{RF03.04} (programación mediante cron), \textbf{RF04.02} (contraste teórico vs. real) y \textbf{RF04.03} (verificación de \textit{items} no continuables).

La iteración resultó más ligera de lo previsto en términos de desarrollo, ya que la mayor parte del esfuerzo se concentró en analizar las tarifas reales de IONOS Cloud y en justificar la elección de tolerancias diferenciadas para \textit{input} y \textit{output}. Una vez tomada esa decisión, la integración fue prácticamente inmediata: la rama \texttt{non-consumable} del \textit{dispatcher} ya estaba disponible y solo fue necesario configurar los SKUs en CouchDB. El testeo, en cambio, se vio condicionado por el calendario mensual de las facturas, pues hubo que esperar al cierre del período para poder ejecutar los escenarios completos sobre datos reales.

No se requirió reunión con el cliente: los cambios se comunicaron de forma asíncrona y la única decisión relevante --- los valores concretos de tolerancia --- se cerró por mensaje, basándose en los datos observados durante la iteración 4.

\begin{table}[H]
    \centering
    \begin{tabular}{|l|c|c|c|}
        \hline
        \textbf{Tarea} & \textbf{Estimadas (h)} & \textbf{Reales (h)} & \textbf{Desviación} \\
        \hline
        Análisis de tarifas y tolerancias       & 4h  & 4h y 30' & +12.5\% \\
        Configuración en CouchDB                & 1h  & 30'      & -50\%   \\
        Integración del tipo \texttt{non-consumable} & 6h  & 4h       & -33.3\% \\
        Comprobación contrato completo          & 1h  & 10'      & -83\%   \\
        Configuración de alertas y rutinas      & 2h  & 2h       & 0\%     \\
        Testeo de funcionamiento                & 6h  & 7h       & +16.7\% \\
        \hline
        \textbf{Total} & \textbf{20h} & \textbf{18h y 10'} & \textbf{-9.2\%} \\
        \hline
    \end{tabular}
    \caption{Comparativa de horas estimadas vs reales en la Iteración 5}
    \label{tab:horas-iteracion5}
\end{table}